\documentclass[12pt,a4paper]{article}
\usepackage{fontspec}
\usepackage{xeCJK}
\usepackage[margin=20mm]{geometry}
\usepackage{xcolor}
\usepackage{array}
\usepackage{colortbl}
\usepackage{fancyhdr}
\setmainfont{Arial}
\setCJKmainfont{SimHei}
\pagestyle{fancy}\fancyhf{}\rhead{设备维护中心}\lhead{维修工单}\cfoot{\thepage}
\setlength{\parindent}{0pt}\setlength{\tabcolsep}{2pt}
\begin{document}
{\fontsize{22}{27}\selectfont\bfseries 设备维修工单}\par\vspace{3mm}
{\color{gray}\small 华东制造中心 / 资产维护系统}\par\vspace{5mm}{\color{black}\rule{\textwidth}{1.2pt}}
\vspace{6mm}
\begin{tabular}{|>{\bfseries}p{28mm}|p{50mm}|>{\bfseries}p{28mm}|p{50mm}|}\hline
\rowcolor{blue!12} \bfseries 字段 & \bfseries 内容 & \bfseries 字段 & \bfseries 内容\\\hline
工单编号 & WO-2026-0817-042 & 工单状态 & 已完成\\\hline
资产编号 & SYNTH-042 & 设备型号 & AX-200 温控模块\\\hline
所属区域 & 二号车间 / B-12 & 报修时间 & 2026-08-17 08:35\\\hline
报修人 & 李明 & 维修负责人 & 周海\\\hline
\end{tabular}
\vspace{8mm}
\section*{一、故障描述}
设备在夜班运行期间出现间歇性过热告警，控制面板在 04:10 和 04:26 两次进入保护模式。现场未发现外壳破损或液体泄漏，生产线已切换到备用模块。
\section*{二、现场初检}
\begin{tabular}{|p{40mm}|p{38mm}|p{40mm}|p{30mm}|}\hline
\rowcolor{gray!15} 检查项目 & 读数 & 标准范围 & 结论\\\hline
环境温度 & 28.4 C & 15--35 C & 正常\\\hline
输入电压 & 23.8 V & 22--26 V & 正常\\\hline
模块温度 & 76.2 C & < 70 C & 异常\\\hline
风扇转速 & 1,120 rpm & > 1,500 rpm & 异常\\\hline
报警代码 & THR-07 & 无 & 已记录\\\hline
\end{tabular}
\newpage
{\fontsize{18}{22}\selectfont\bfseries 设备维修工单 / 检查记录}\par\vspace{8mm}
\section*{三、诊断过程}
\begin{enumerate}
\item 断开主电源并等待 90 秒，确认直流母线电压降至 0 V；
\item 拆下右侧风扇护罩，检查叶片、轴承和排风通道；
\item 使用 FLUKE-87V 测量输入端、控制端和热敏电阻阻值；
\item 读取控制器日志，确认 THR-07 与风扇低速状态同时出现；
\item 将备用风扇接入测试端口，观察温度曲线 15 分钟。
\end{enumerate}
\section*{四、测量记录}
\begin{tabular}{|p{37mm}|p{34mm}|p{34mm}|p{42mm}|}\hline
\rowcolor{gray!15} 时间 & 模块温度 & 风扇转速 & 备注\\\hline
09:02 & 75.8 C & 1,105 rpm & 原风扇，持续报警\\\hline
09:08 & 54.1 C & 1,860 rpm & 接入备用风扇\\\hline
09:16 & 55.0 C & 1,842 rpm & 负载 60\%\\\hline
09:24 & 61.7 C & 1,805 rpm & 负载 85\%\\\hline
09:31 & 63.2 C & 1,792 rpm & 无报警\\\hline
\end{tabular}
\section*{五、处理措施}
原风扇型号 FAN-AX200-12V 的轴承阻力明显升高，已更换为同规格备件。清理排风口积尘后重新固定护罩，恢复控制器参数并执行满载测试。
\newpage
{\fontsize{18}{22}\selectfont\bfseries 设备维修工单 / 验收与签字}\par\vspace{8mm}
\section*{六、更换零件}
\begin{tabular}{|p{40mm}|p{38mm}|p{28mm}|p{38mm}|}\hline
\rowcolor{gray!15} 零件名称 & 规格/型号 & 数量 & 追踪编号\\\hline
冷却风扇 & FAN-AX200-12V & 1 & SP-2026-0817-11\\\hline
固定螺钉 & M4 x 12 不锈钢 & 4 & SP-2026-0817-12\\\hline
\end{tabular}
\section*{七、验收结果}
\begin{tabular}{|p{58mm}|p{56mm}|p{28mm}|}\hline
\rowcolor{gray!15} 验收项目 & 结果 & 状态\\\hline
空载运行 10 分钟 & 温度 49.6 C，无报警 & 通过\\\hline
额定负载运行 30 分钟 & 温度 63.2 C，转速 1,792 rpm & 通过\\\hline
断电恢复测试 & 参数保持，日志无新增错误 & 通过\\\hline
外观与防护 & 护罩闭合，接地线牢固 & 通过\\\hline
\end{tabular}
\section*{八、结论}
故障原因为冷却风扇轴承磨损导致转速不足。更换风扇后设备恢复正常，建议在 2026-09-01 前将同批次模块纳入预防性巡检。
\vspace{8mm}\par\noindent
\begin{tabular}{p{50mm}p{50mm}p{50mm}}
维修负责人：周海 & 生产确认：李明 & 复核：王倩\\[12mm]
日期：2026-08-17 & 日期：2026-08-17 & 日期：2026-08-17
\end{tabular}
\end{document}